% GENERATED FILE. Edit canonical lessons, then run npm run generate:books.
% canonical-source-hash: fnv1a64:1c5d7ddf3229d9bc
% canonical-lessons: SA-C219-bilvam, SA-C219-narangam, SA-C219-kapitthah, SA-C219-sevaphalam, SA-C219-cipitakah

\chapter{Bael fruit, Orange, Wood apple, Apple, Flattened rice}
\label{ch:sa-bael-fruit-orange-wood-apple-apple-flattened-rice}
\hlchaptermodality{219}

\begin{chapteropening}
\textbf{By the end of this chapter:} \emph{I can say a bael fruit, an orange, a wood apple, an apple and flattened rice.}

Complete the last lesson of chapter 219: I can say a bael fruit, an orange, a wood apple, an apple and flattened rice.
\end{chapteropening}

\section[\texorpdfstring{\sk{बिल्वम्} (bilvam)}{बिल्वम् (bilvam)}]{\sk{बिल्वम्} (bilvam) — a bael fruit}

\label{lesson:SA-C219-bilvam}



\begin{quote}
Before the new one: say the Sanskrit for a plough, then the Sanskrit for a sickle.
\end{quote}

\subsection*{You'll want to know: \sk{बिल्वम्}}
\textbf{\sk{बिल्वम्}} — \emph{bilvam} — \textquotedblleft{}a bael fruit\textquotedblright{}.

\begin{grammarlens}[title={What you've built}]
One more word for this chapter.
\end{grammarlens}

\subsection*{Your turn}
\begin{itemize}
\raggedright
  \item \emph{Say it:} \emph{bilvam}
  \item \emph{Say it:} \emph{bilvam}, once more
  \item \emph{Say it:} say \emph{bilvam}
  \item \emph{Recall:} say the Sanskrit for a plough, then the Sanskrit for a sickle, then say \emph{bilvam} again
\end{itemize}

\begin{tcolorbox}[breakable,colback=teal!4,colframe=teal!35!black,title={Before you move on}]
What does \sk{बिल्वम्} mean? (\textquotedblleft{}A bael fruit\textquotedblright{}.) Say it once more.
\end{tcolorbox}
\section[\texorpdfstring{\sk{नारंगम्} (nāraṅgam)}{नारंगम् (nāraṅgam)}]{\sk{नारंगम्} (nāraṅgam) — an orange}

\label{lesson:SA-C219-narangam}



\begin{quote}
Before the new one: say the Sanskrit for a sickle, then the Sanskrit for a bael fruit.
\end{quote}

\subsection*{You'll want to know: \sk{नारंगम्}}
\textbf{\sk{नारंगम्}} — \emph{nāraṅgam} — \textquotedblleft{}an orange\textquotedblright{}.

\begin{grammarlens}[title={What you've built}]
One more word for this chapter.
\end{grammarlens}

\subsection*{Your turn}
\begin{itemize}
\raggedright
  \item \emph{Say it:} \emph{nāraṅgam}
  \item \emph{Say it:} \emph{nāraṅgam}, once more
  \item \emph{Say it:} say \emph{nāraṅgam}
  \item \emph{Recall:} say the Sanskrit for a sickle, then the Sanskrit for a bael fruit, then say \emph{nāraṅgam} again
\end{itemize}

\begin{tcolorbox}[breakable,colback=teal!4,colframe=teal!35!black,title={Before you move on}]
What does \sk{नारंगम्} mean? (\textquotedblleft{}An orange\textquotedblright{}.) Say it once more.
\end{tcolorbox}
\section[\texorpdfstring{\sk{कपित्थः} (kapitthaḥ)}{कपित्थः (kapitthaḥ)}]{\sk{कपित्थः} (kapitthaḥ) — a wood apple}

\label{lesson:SA-C219-kapitthah}



\begin{quote}
Before the new one: say the Sanskrit for a bael fruit, then the Sanskrit for an orange.
\end{quote}

\subsection*{You'll want to know: \sk{कपित्थः}}
\textbf{\sk{कपित्थः}} — \emph{kapitthaḥ} — \textquotedblleft{}a wood apple\textquotedblright{}.

\begin{grammarlens}[title={What you've built}]
One more word for this chapter.
\end{grammarlens}

\subsection*{Your turn}
\begin{itemize}
\raggedright
  \item \emph{Say it:} \emph{kapitthaḥ}
  \item \emph{Say it:} \emph{kapitthaḥ}, once more
  \item \emph{Say it:} say \emph{kapitthaḥ}
  \item \emph{Recall:} say the Sanskrit for a bael fruit, then the Sanskrit for an orange, then say \emph{kapitthaḥ} again
\end{itemize}

\begin{tcolorbox}[breakable,colback=teal!4,colframe=teal!35!black,title={Before you move on}]
What does \sk{कपित्थः} mean? (\textquotedblleft{}A wood apple\textquotedblright{}.) Say it once more.
\end{tcolorbox}
\section[\texorpdfstring{\sk{सेवफलम्} (sevaphalam)}{सेवफलम् (sevaphalam)}]{\sk{सेवफलम्} (sevaphalam) — an apple}

\label{lesson:SA-C219-sevaphalam}



\begin{quote}
Before the new one: say the Sanskrit for an orange, then the Sanskrit for a wood apple.
\end{quote}

\subsection*{You'll want to know: \sk{सेवफलम्}}
\textbf{\sk{सेवफलम्}} — \emph{sevaphalam} — \textquotedblleft{}an apple\textquotedblright{}.

\begin{grammarlens}[title={What you've built}]
One more word for this chapter.
\end{grammarlens}

\subsection*{Your turn}
\begin{itemize}
\raggedright
  \item \emph{Say it:} \emph{sevaphalam}
  \item \emph{Say it:} \emph{sevaphalam}, once more
  \item \emph{Say it:} say \emph{sevaphalam}
  \item \emph{Recall:} say the Sanskrit for an orange, then the Sanskrit for a wood apple, then say \emph{sevaphalam} again
\end{itemize}

\begin{tcolorbox}[breakable,colback=teal!4,colframe=teal!35!black,title={Before you move on}]
What does \sk{सेवफलम्} mean? (\textquotedblleft{}An apple\textquotedblright{}.) Say it once more.
\end{tcolorbox}
\section[\texorpdfstring{\sk{चिपिटकः} (cipiṭakaḥ)}{चिपिटकः (cipiṭakaḥ)}]{\sk{चिपिटकः} (cipiṭakaḥ) — flattened rice}

\label{lesson:SA-C219-cipitakah}



\begin{quote}
Before the new one: say the Sanskrit for a wood apple, then the Sanskrit for an apple.
\end{quote}

\subsection*{You'll want to know: \sk{चिपिटकः}}
\textbf{\sk{चिपिटकः}} — \emph{cipiṭakaḥ} — \textquotedblleft{}flattened rice\textquotedblright{}.

\begin{grammarlens}[title={What you've built}]
One more word for this chapter.
\end{grammarlens}

\subsection*{Your turn}
\begin{itemize}
\raggedright
  \item \emph{Say it:} \emph{cipiṭakaḥ}
  \item \emph{Say it:} \emph{cipiṭakaḥ}, once more
  \item \emph{Say it:} say \emph{cipiṭakaḥ}
  \item \emph{Recall:} say the Sanskrit for a wood apple, then the Sanskrit for an apple, then say \emph{cipiṭakaḥ} again
\end{itemize}

\begin{tcolorbox}[breakable,colback=teal!4,colframe=teal!35!black,title={Before you move on}]
What does \sk{चिपिटकः} mean? (\textquotedblleft{}Flattened rice\textquotedblright{}.) Say it once more.
\end{tcolorbox}
